%%%%%%%%%%%%%%%%%%%%%%%%%%%%%%%%%%%%%%%%%%%%%%%%%%%%%%%%%%%%
% Task 02 : Single View Geometry and Distance Estimation
%%%%%%%%%%%%%%%%%%%%%%%%%%%%%%%%%%%%%%%%%%%%%%%%%%%%%%%%%%%%

\section{Overview}

In this task, you will estimate real-world distances using a single
image by -- well -- applying what you learnt in class. The task
focuses on  linearity, cross-ratio, and the practical use
of object localization (using a  deep learning tool) to infer
distances from a single image. This topic is called single view metrology.

\section{Background}

You are provided with an image (see Fig.~\ref{fig:bicycleBBox}) of
the way from the main building
(MB) of IIT Bombay to the arch of knowledge (gyanam-paramam-dhyeyam) (``arch''). (As an aside, the arch encodes the motto of IITB).

Take a moment to digest the provided image. Observe
\begin{itemize} 
\item the red line overlaid for reference, 
\item the cross road (``intersection'') that has the palm tree on one side and a building
on the other side
\item a bicycle
\end{itemize}

Amongst the  objects, look for the street lights  towards the right. We refer to these as 
``poles'' in this assignment. The poles are numbered starting from the camera to the Arch. We will casually refer to the red line as equivalent to MB.
We are
going to assume that the cross road (aka intersection)  is a straight line passing
through the base of the third pole parallel to the red
line, and that the mb-to-arch line is perpendicular to
the intersection line.


%%%%%%%%%%%%%%%%%%%%%%%%%%%%%%%%%%%%%%%%%%%%%%%%%%%%%%%%%%%%
\section*{Honour code and sources}
I followed the honour code. The numbers, the measurements, the code and
the writing here are my own work. Sources I leaned on: the course
lectures of 03-Jun and 10-Jun (camera matrix, linearity, cross ratio)
and the 25-May surveillance-camera worked example; Ultralytics YOLOv8
documentation (\url{https://docs.ultralytics.com/tasks/detect}); the
COCO class list (\url{https://cocodataset.org}); Google Maps for the
ground-truth distance; and MS Paint / Python (Pillow) for reading pixel
coordinates and drawing the annotation. I also used an LLM as a sounding
board to clarify the cross-ratio derivation and to debug LaTeX package
issues; I did not let it stand in for doing the geometry or the
measurements, which I worked through myself and cross-checked two ways.

%%%%%%%%%%%%%%%%%%%%%%%%%%%%%%%%%%%%%%%%%%%%%%%%%%%%%%%%%%%%
\section{Who \href{https://en.wikipedia.org/wiki/Who_Moved_My_Cheese\%3F}{moved} my bicycle?}%

\subsection{Task}
Consider the bicycle in
Fig.~\ref{fig:bicycleBBox}. Report the real-world distance between the
bicycle and the arch.

\begin{figure}[h!]
\centering
\includegraphics[width=0.45\linewidth]{bicycle}
\caption{A bounding box prediction from YOLOv8.}
\label{fig:bicycleBBox}
\end{figure}

Use a pretrained deep learning model to detect the
bicycle. This is an object detection task. Please refer to tutorial here, and stick to \textbf{YOLO} models: \url{https://docs.ultralytics.com/tasks/detect#models}
Also look at the COCO dataset \url{https://cocodataset.org/#home} to find out the class identifiers for bicycles and persons.
%The bounding box from the model will give you the coordinates of the bicycle. 

\begin{studentSpace}
\noindent\textbf{Answer.}
My estimate is that the bicycle is about \textbf{120 metres} from the
arch (more precisely 121--124 m by two independent calculations; see
below). The short version of how I got there: detect the bicycle with
YOLO, take the point where its wheels touch the ground, and use the
cross ratio along the MB--arch line to convert that pixel position into
a real distance.

The image of the MB--arch line is the central axis of the avenue, and
the four collinear points I work with are, in order of distance from the
camera: $A'$ (the red MB line), $B'$ (the bicycle), $C'$ (the
intersection at the third pole) and $V'$ (the vanishing point, which
here sits right at the arch). I drew this line in dark blue and marked
the four points -- see Fig.~\ref{fig:blueline}.

\begin{figure}[h!]
\centering
\includegraphics[width=0.55\linewidth]{otherImage1}
\caption{The dark-blue MB--arch axis with the four collinear points
$A'$ (MB), $B'$ (bicycle), $C'$ (intersection), $V'$ (vanishing point).
Pixel rows: $A'=3500$, $B'=2150$, $C'=2050$, $V'=1720$ (image is
$2268\times4032$). The magenta box is the YOLO detection.}
\label{fig:blueline}
\end{figure}

The cross ratio of four collinear points is invariant to projective
transformation (Lecture~05, 25-May). Using the definition given in class,
\[
C(x_1,x_2,x_3,x_4)\;\equiv\;\frac{|x_1x_2|\,|x_3x_4|}{|x_1x_3|\,|x_2x_4|},
\]
and applying it to $(x_1,x_2,x_3,x_4)=(A',B',C',V')$ exactly as in the
surveillance-camera example, the image cross ratio read off the pixel
rows is
\[
C=\frac{A'B'\cdot C'D'}{A'C'\cdot B'D'}
 =\frac{|y_{A'}-y_{B'}|\,|y_{C'}-y_{V'}|}{|y_{A'}-y_{C'}|\,|y_{B'}-y_{V'}|}
 =\frac{1350\cdot 330}{1450\cdot 430}=0.71,
\]
where $D'\equiv V'$ is the vanishing point. In the real world $V'$ is the
point at infinity, so (as done in class) the two terms containing it
cancel, and the cross ratio reduces to a ratio of two known distances:
\[
C=\frac{A B\cdot C D}{A C\cdot B D}
 =\frac{s_B\cdot\infty}{d_{\text{int}}\cdot\infty}
 =\frac{s_B}{d_{\text{int}}}
\;\;\Longrightarrow\;\;
s_B=d_{\text{int}}\cdot C=50\times0.71=35.7~\text{m},
\]
where $s_B$ is the bicycle's distance from MB and $d_{\text{int}}=50$~m
is the MB--intersection reference. Hence the bicycle is
$160-35.7\approx\mathbf{124}$~m from the arch. As a check I re-did the
same cross-ratio equation with the vanishing point replaced by the arch
(a finite point at $160$~m), i.e.\ MB, intersection and arch as three
finite references; that gives $\approx 121$~m. The two agree to within a
few metres, so I report roughly $120$~m.
\end{studentSpace}


\subsection{Questions}

You are expected to write answers inline in the student space provided. (see Section 5)

\begin{enumerate}
\item Use a tool from Google that tells you the real world distance from MB (the red line) to the arch.  Use a number in metres and approximate to the nearest distance divisible by 10. 
\begin{studentSpace}
\noindent I used Google Maps in satellite view with the ``Measure
distance'' tool (right-click on the start point, then click the end
point). Measuring from the Main Building red line straight down the
tree-lined avenue to the arch reads \textbf{160.35~m}, which rounds to
\textbf{160~m}. The screenshot is in Fig.~\ref{fig:map}.
\begin{figure}[h!]
\centering
\includegraphics[width=0.85\linewidth]{otherImage2}
\caption{Google Maps ``Measure distance'': MB (0~m marker) to the arch
end of the avenue $=160.35$~m $\approx 160$~m.}
\label{fig:map}
\end{figure}
\end{studentSpace}

\item Why can you NOT find the distance to the bicycle without making any assumption about the bicycle using the same tool? (Answer in one line.)
\begin{studentSpace}
\noindent Because Google's tool only measures between fixed, mapped
locations, and the bicycle is a movable object with no map coordinates,
so without assuming where it lies (e.g.\ that it sits on the MB--arch
line) there is nothing for the tool to measure to.
\end{studentSpace}

\item Ensure that you draw a line in dark blue color (using an image annotation tool of your choice) that is used for obtaining the cross reference.  Show this line and the corresponding image in your submission. How did you draw the line? Keep in mind that cross-ratio is a \textbf{1D} concept and requires 4 collinear points. Make sure these points lie on the dark blue color line above.
\begin{studentSpace}
\noindent The dark-blue line (Fig.~\ref{fig:blueline}) is the image of
the MB--arch axis -- the central line of the avenue. I drew it like this:
first I extended the two converging edges of the path (the hedge line on
the left and the curb/flower-border on the right) until they met; that
meeting point is the vanishing point $V'$, sitting almost exactly under
the arch. I then drew one straight line from where that axis crosses the
red MB line, up through the bicycle's wheel-contact point, to $V'$.
Because the cross ratio is a one-dimensional invariant, all four points
have to be collinear; placing $A',B',C',V'$ on this single line
guarantees that, and it also means only each point's position along the
line matters. I read the points off in MS Paint / with Pillow and used
the pixel \emph{row} ($y$) of each, since the cross ratio is unchanged by
any affine re-scaling of the line.
\end{studentSpace}

\item Ensure that you annotate 4 points of interest as A', B', C', V'  in sorted order of distance from the camera where B' corresponds to the biycle and V' corresponds to the vanishing point.  What do A' and B' correspond to?
We recommend using MS Paint for annotation and identifying pixel coordinates.
\begin{studentSpace}
\noindent Sorted by distance from the camera (nearest first):
$A'$ = the MB / red line (the near reference, $0$~m);
$B'$ = the bicycle (its wheel/ground-contact point, the unknown);
$C'$ = the intersection at the third-pole base ($50$~m reference);
$V'$ = the vanishing point (image of the point at infinity along the
avenue, coinciding with the arch). So directly: \textbf{$A'$ is the MB /
red-line reference point (start of measurement, 0~m), and $B'$ is the
bicycle's wheel/ground-contact point.}
\end{studentSpace}

\item Be sure to  explain in detail your method (and code) for finding the bicycle. Does the choice of reference point within the bounding box from the deep learning model affect distance estimation accuracy? Quantify.
\end{enumerate}
\begin{studentSpace}
\noindent\textbf{Finding the bicycle.} I ran a pretrained YOLOv8 detector
(Ultralytics) on the unsegmented image. YOLOv8 is trained on COCO, where
\texttt{bicycle} is class id $1$ (and \texttt{person} is $0$), so I keep
only that class. The detector returns a bounding box
$\texttt{(x1,y1,x2,y2)}=(1105,2013,1204,2153)$ with confidence $0.71$.
From this box I take the \emph{bottom-centre} point $(1154,2153)$ as the
reference, because that is where the bicycle touches the ground plane on
which the cross-ratio line lives. That point becomes $B'$, and I feed its
row into the cross-ratio formula above. The two short programs are
\texttt{code/detect\_bicycle.py} (detection $\rightarrow$ box
$\rightarrow$ ground-contact point) and
\texttt{code/crossratio\_distance.py} (cross-ratio solve and the
sensitivity sweep).

\noindent\textbf{Does the reference point inside the box matter? Yes,
a lot.} Keeping everything else fixed and only moving the reference point
across the $140$-pixel-tall box:

\begin{center}
\begin{tabular}{lcc}
\hline
reference point & pixel row $y_{B'}$ & bicycle $\rightarrow$ arch \\
\hline
top (saddle)          & 2010 & 101.5 m \\
box centre            & 2080 & 115.1 m \\
wheel / ground contact& 2150 & 124.3 m \\
\hline
\end{tabular}
\end{center}

That is a spread of \textbf{22.7~m over 140 px}, about $0.16$~m per pixel
of vertical error. The sensitivity is large because near the vanishing
point the image is heavily compressed, so a small pixel shift maps to a
big real-world shift. The wheel/ground-contact point is the correct
choice: the metrology line lives on the ground, so a point lifted off the
ground (saddle, box centre) reads as ``higher'' in the image, which
biases the bicycle nearer and inflates the remaining distance to the
arch. As a sanity note, rounding all four pixel rows to the nearest 10
changed the final distance by less than $1$~m, so the dominant error is
the choice of reference point, not pixel-reading precision.
\end{studentSpace}


%%%%%%%%%%%%%%%%%%%%%%
\section{Your picture}

In this section, take a picture of you and, alternately, your teammate
(or some other person) similar to the one that has been provided in
your locality. Identify two specific points of interest and tell us
how far s/he is from the point of interest in the real world. We are
expecting two pictures in your submission; one corresponding to the
problem statement (often called as the first person view), and the
second, a picture of you taking a picture of your teammate, i.e., a
third person view. The second picture will show both the cameraman and
the teammate (and ideally at the same time).

\subsection{Tasks}

\begin{enumerate}
\item Capture and include the images in your submission.
\item Detect the teammate using the same deep learning model.
\item Estimate the distance between one of the objects of interest and
  the team mate.
\begin{studentSpace}
\noindent The first-person view is Fig.~\ref{fig:fpv} and the
third-person view is Fig.~\ref{fig:tpv}. The receding line is the central
walkway outside building B-6 in my society. My two points of interest on
that line are the \textbf{camera position} where I stood (the near point
$A'$, taken as $0$~m) and the \textbf{far end of the walkway} where it
meets the distant towers past the footbridge (the vanishing region,
$V'$), which is roughly $100$~m away. I detected the person with the
same detector (\texttt{detect\_bicycle.py}, kept class switched to
\texttt{person}, COCO id $0$); the bottom-centre of that box is the
ground-contact point $B'\approx(1540,1730)$. (I worked on this assignment individually; the person in the photos is a family member who simply stood in.) Drawing the dark-blue line
through $A'$, $B'$ and the vanishing point, and applying the cross ratio
with $V'$ as the point at infinity and the $\sim 100$~m depth as the
scale, the person comes out at roughly $\mathbf{7\text{--}8}$~m from the
camera. That agrees with the third-person view (Fig.~\ref{fig:tpv}),
where the person and I are standing about $7$--$8$~m apart -- an
independent check that the single-image estimate is right.

\begin{figure}[h!]
\centering
\includegraphics[width=0.55\linewidth]{firstPersonView}
\caption{First-person view in my locality, with the dark-blue line and
the points $A'$ (near, 0~m), $B'$ (the person) and $V'$ (vanishing point,
$\sim 100$~m) marked (taken at dusk, $\sim$7:25 PM).}
\label{fig:fpv}
\end{figure}
\begin{figure}[h!]
\centering
\includegraphics[width=0.55\linewidth]{thirdPersonView}
\caption{Third-person view: me photographing the other person; we are about 7--8 m apart.}
\label{fig:tpv}
\end{figure}
\end{studentSpace}
\end{enumerate}

You are expected to write answers inline.

\subsection{Question}
\begin{enumerate}

\item As before, draw the line of interest (in blue color) marking all points of interest. 
  
\item Explain how you got the real world distances.
\begin{studentSpace}
\noindent The one real distance I fixed is the depth of the walkway to
the vanishing region -- about $100$~m, which I estimated by pacing along
the avenue (a private society driveway is not easy to measure on Google
Maps, unlike the open campus avenue in Part~1). The person's distance
is \emph{not} measured directly -- it is computed: the near point and the
vanishing point fix the projective frame on the receding line, and the
cross ratio converts the person's ground-contact pixel position into
metres, giving $\approx 7$--$8$~m. I then checked that against the
third-person photo, where the two of us stand about $7$--$8$~m apart. So
exactly one distance is assumed and everything else is geometry that the
second photograph independently confirms.
\end{studentSpace}

\item Discuss practical challenges encountered while applying the projective
geometry concepts to real-world images. Number all challenges and write independent challenges i.e. don't mix up the challenges.
\begin{studentSpace}
\begin{enumerate}
\item \textbf{Locating the vanishing point.} Real path edges are not
perfectly straight (hedges bulge, curbs wander), so the two edges did not
cross at a single crisp point. I had to fit the straighter middle
sections and accept a small band of uncertainty for $V'$.
\item \textbf{Choosing the reference point on the detected object.} The
YOLO box is a rectangle; the object touches the ground only at the
bottom. Picking the box centre instead of the wheel/feet shifted the
answer by tens of metres (quantified in Q5).
\item \textbf{Perspective compression near the horizon.} Close to the
vanishing point a single pixel covers many metres, so any object far down
the line is intrinsically hard to localise, and small reading errors blow
up.
\item \textbf{The fronto-parallel / on-the-line assumption.} The method
assumes the object sits exactly on the chosen line and the camera has no
roll. If the bicycle (or teammate) is even slightly off the avenue axis,
or the phone is tilted, the collinearity that the cross ratio depends on
is violated.
\item \textbf{Getting a trustworthy real-world reference.} The whole
scale rides on one looked-up distance (160~m here). A measurement error
or a different rounding in Google Maps scales the entire result
proportionally.
\end{enumerate}
\end{studentSpace}

\item What programs did you write and why?  
\begin{studentSpace}
\noindent Two small Python programs, both in \texttt{code/}.
\texttt{detect\_bicycle.py} loads pretrained YOLOv8 and returns the
bicycle (or person) bounding box and its bottom-centre ground-contact
point -- I wrote it so the object localisation is reproducible rather
than hand-clicked, and so the same script works for Part~1 and Part~2 by
swapping the COCO class. \texttt{crossratio\_distance.py} takes the four
pixel rows and the reference distances, computes the cross ratio, solves
for the unknown distance two ways (with and without the vanishing point),
and runs the bounding-box sensitivity sweep used in Q5. I split them so
detection (heavy, needs the model) and geometry (instant, pure
arithmetic) stay independent. I deliberately did not bundle the model
weights; they download from the internet on first run.
\end{studentSpace}

\end{enumerate}

\section{Submit}

Write a description of your method and all relevant aspects of your
experiments. Write it in line in the text under the relevant section.  Your writing should be similar to  a blog post so that anyone in the class can fully understand the pain/joy/content you experienced (feel free to take copious input from an LLM but do not substitute an LLM for learning).

The total length should be no less than 500 words and
should normally exceed 800 words.  Your ultimate aim in writing your answers
to mention important aspect of your assignment for the task, with
emphasis on the last task.  This will include the methods employed
listing all the software tools you employed, as well experimental
aspects such as when you took the picture, how many pictures 
you took, whether the method works in dawn, morning, and dusk, and so
on. Be creative.  

Do not submit any large file such as
model weights and ancillary files that we can easily get from the
internet to reproduce.

\begin{studentSpace}
\noindent\textbf{A short blog: measuring a campus avenue from one photo.}

When I first read the title (``Who moved my bicycle?'') I thought the
hard part would be the deep learning. It wasn't. Detecting the bicycle
took one call to a pretrained model. The real work, and the fun, was in
the geometry: convincing myself that you can recover a real distance in
metres from a flat photograph, with no depth sensor, no second camera,
just one image and one known length.

The reason it works is the cross ratio. In a photograph, perspective
ruins ordinary measurement -- equal real distances do not stay equal on
the image, which is obvious the moment you look at how the avenue's edges
squeeze together towards the arch. So ``count the pixels and scale'' is
hopeless. But there is exactly one quantity built from four points on a
line that survives projection unchanged: the cross ratio. That is the
whole trick. If I can find four collinear points and I know three of the
real distances between them, the fourth falls out. This is the same idea
as the surveillance-camera example from the 25-May class, where the train
gets located between two stations using the point where the rails
converge.

My four points, in order away from the camera, were $A'$ the Main
Building red line, $B'$ the bicycle, $C'$ the intersection at the third
pole, and $V'$ the vanishing point -- which, conveniently, sits almost
exactly at the arch. The first thing I learned the hard way is that the
arch and the vanishing point are essentially the same point in this
photo, so I could not use ``MB, arch, vanishing point'' as my three
references; the arch and $V'$ were too close in pixels and the arithmetic
became unstable. Switching to the intersection (well separated from the
arch in the image) as the middle reference fixed that. A second neat
realisation: the cross ratio doesn't care how you put coordinates on the
line, so I only needed the pixel \emph{row} of each point, not both
coordinates. That cut the measurement effort in half and made the dark
blue annotation line easy to draw and read.

Tools. For detection I used Ultralytics YOLOv8, with the COCO class ids
(bicycle = 1, person = 0). For pixel coordinates I used MS Paint and a
short Pillow script; for the ground-truth distance, Google Maps' measure
tool gave 160.35~m from MB to the arch, which I rounded to 160~m. The
geometry itself is a few lines of Python. I cross-checked the bicycle
result two ways -- once using the vanishing point as the point at
infinity (124~m), once using three finite landmarks with no vanishing
point (121~m) -- and they agreed to within a few metres, which gave me
confidence the method, and my reading of the points, were sound. I report
the bicycle as roughly 120~m from the arch.

The part that surprised me most was how much the \emph{choice of point on
the detected object} matters. The YOLO box is a rectangle, but the
bicycle only touches the ground at the bottom. When I tried the box
centre or the saddle instead of the wheel contact, the answer moved by
more than twenty metres. That is not a rounding issue -- it is geometry:
near the vanishing point the image is so compressed that one pixel is
worth about 0.16~m, so where exactly you ``stand'' the object changes a
lot. The lesson is that the reference point must be the ground-contact
point, every time. By contrast, rounding the pixel rows to the nearest 10
barely moved the result, so the human pixel-reading was not the
bottleneck; the modelling choice was.

For the second part I repeated the whole pipeline on my own photographs
in my housing society, on the walkway outside building B-6. I took a
first-person view with another person standing on the receding walkway, and
a third-person view showing me in the act of photographing them so both
of us are in frame. Both shots were taken at dusk, around 7:25 PM, and
that mattered: in the fading light the far end of the walkway and my
person's feet were noticeably harder to pin down than the bright
daytime campus shot in Part~1, which widened my error on the vanishing
point. Taking the depth of the walkway as about 100~m, the cross ratio
placed the person roughly 7--8~m from the camera, and the third-person
photo -- where the two of us stand about 7--8~m apart -- agreed, which
was satisfying to see. The same two scripts handled detection and
distance; the only change was keeping the \texttt{person} class instead
of \texttt{bicycle}.

If I had to summarise the experience: the deep learning is a convenience,
not the substance. The substance is one old, elegant projective
invariant, and the discipline of choosing the right point and the right
references. Most of my errors came not from the maths but from the
assumptions underneath it -- a not-quite-straight path edge, an object
slightly off the line, a tilted phone. Respect those assumptions and a
single photo really does give you metres.
\end{studentSpace}
